% !TEX root = ../main.tex
\chapter{Corrispondenza di Curry-Howard}
\label{chap:curry_howard_isomorphism}

\section{Definizione generale di isomorfismo}
Il concetto di \emph{isomorfismo} è uno dei pilastri della matemantica. Esso identifica una relazione di equivalenza 
profonda tra due strutture che, pur operando in domini distinti e/o utilizzando notazioni distinte, condividono la 
medesima architettura interna. In questa sezione definiremo le condizioni necessarie e sufficenti affinche due sistemi 
formali possano considerarsi isomorfi. 

\subsection{Biunivocità e preservazione della struttura}
In termini algebrici classici, un isomorfismo tra due strutture di $\mathcal{X}$ e $\mathcal{Y}$ è definito come un 
\emph{omomorfismo biunivoco}. Se consideriamo due insiemi dotati di operazioni interne $(\mathcal{X}, C_{\mathcal{X}})$ 
e $(\mathcal{Y}, C_{\mathcal{Y}})$, un funzione $f: \mathcal{X} \to \mathcal{Y}$ è un isomorfismo se soddisfa:

\begin{description}[leftmargin=1.5em, style=nextline, font=\bfseries]
    \item[Biunivocità] 
        La funzione $f$ deve essere una biiezione. Questo garantisce l'esistenza di una funzione inversa $f^{-1}$, permettendo
        un processo di \emph{round-trip} tra due domini. In un sistema formale, ciò assicura che non vi sia perdita di contenuto 
        informativo, ogni elemento del primo sistema ha un'immagine unica nel secondo e viceversa.
    \item[Omomorfismo]
        La funzione deve preservare la struttura delle operazioni. Per n-upla di elementi $x_1, \dots, x_n \in \mathcal{X}$ deve 
        valere l'identita, $f(C_{\mathcal{X}}(x_1, \dots, x_n)) = C_{\mathcal{Y}}(f(x_1), \dots, f(x_n))$.
\end{description}

\section{Introduzione isomorfismo di Curry-Howard}
L'isomorfismo di Curry-Howard stabilisce un'identità tra la teoria della dimostrazione e la teoria del tipi. Secondo tale 
principio, una proposizione logica non è una affermazione binaria, ma un \emph{tipo}. Tale identità porta una dimostrazione 
a essere un \emph{termine} appartenente a un determinato tipo. Questa corrispondenza trasforma l'attività logica e quella 
computazionale in due manifestazioni del medesimo processo formale.
\begin{table}[H]
    \centering
    \renewcommand{\arraystretch}{1.3}
    \begin{tabular}{lll}
        \hline
        \textbf{Categoria} & \textbf{Logica}   & \textbf{Informatica} \\ \hline
        Oggetto astratto   & Proposizione $A$  & Tipo $\tau$          \\
        Oggetto concreto   & Dimostrazione $d$ & Termine $M$          \\
        Validità           & $d$ prova $A$     & $M : \tau$           \\
        Assunzioni         & Ipotesi           & Contesto $\Gamma$    \\
        Provabilità        & Teorema           & Abitabilità          \\
        Contraddizione     & Falso ($\bot$)    & Tipo vuoto           \\
        Dinamica           & Semplificazione   & $\beta$-riduzione    \\
        Risultato          & Prova normale     & Forma normale        \\ \hline
    \end{tabular}
    \caption{Corrispondenze sintetiche dell'isomorfismo.}
    \label{tab:curry_howard_isomorphism}
\end{table}\noindent
Questa mappatura implica che la validità di una formula coincida con la costruzione di un oggetto matematico, il termine,
che ne testimonia la correttezza. Il \emph{type checking} agisce quindi come un verificatore meccanico di prove logiche.
L'importanza di tale isomorfismo risiede nella possibilità di traslare i risultati di un dominio nell'altro. Il dominio
della programmaizone riceve rigore logico-matematico, invece il dominio della logica riceve una connotazione algoritmica con il
quale verificare le dimostrazioni.

\subsection{Indipendeza dalla semantica}
Un aspetto fondamentale dell'isomorfismo di Curry-Howard è la sua natura puramente \emph{sintattica}. La corrispondenza non dipende 
dall'interpretazione semantica dei termini, ma esclusivamente dalle regole di formazione e manipolazione dei simboli. In questo contesto, 
la validità di una dimostrazione logica non è legata al concetto di verità assoluta in un modello, bensì alla corretta applicazione delle 
regole di inferenza. Nel $\lambda$-calcolo, la correttezza di un programma è garantita dal rispetto dei vincoli sintattici imposti dal 
sistema di tipi. Questa indipendenza permette di traslare le proprietà del $\lambda$-calcolo tipato in garanzie di coerenza logica. Se un 
sistema di tipi è corretto (sound), esso non permetterà mai la costruzione di un termine che abiti il tipo vuoto ($\bot$), indipendentemente 
dalla semantica operazionale del linguaggio sottostante. L'effetto di maggiore impatto di tale indipendeza è la possibilta di poter utilizzare 
la logica intuizionistica e le sue regole per la dimostrazione del'isomorfismo. Nonostante in questo momento della trattazione su tale logica 
non sia ancora definita una semantica rigorosa. 

\subsection{Metodologia dimostrativa}
La metodologia utilizzata per stabilire l'isomorfismo tra i sistemi di tipo e i sistemi logici si basa sulla comparazione strutturale tra
gli \emph{alberi di derivazione} della deduzione naturale e gli \emph{alberi di tipizzazione} del $\lambda$-calcolo. I type checking nella 
precedente trattazione sono stati definiti come oggetti operazionali, un algortmi ricorsivo. Tale definizione rende implicito l'albero di 
tipizzazione, la struttura della computazione, che pero è il vero oggetto del'isomorfismo. La dimostrazione per induzione sulla struttura 
degli alberi di derivazione avrà come obiettivo la preservazione dei seguenti costrutti e proprietà fondamentali.

\paragraph{Corrispondenze statiche.}
\begin{description}[leftmargin=1.5em, style=nextline, font=\bfseries]
    \item[Isomorfismo dei giudizi] 
    Si stabilisce l'identità formale $\Gamma \vdash M : A \iff \Gamma \vdash A$. Il contesto $\Gamma$ rappresenta simultaneamente l'insieme 
    delle dichiarazioni di tipo per le variabili e l'insieme delle ipotesi logiche attive. In questo senso, il termine $M$ non è un elemento 
    esterno, ma la prova stessa della proposizione $A$ derivata sotto le ipotesi contenute in $\Gamma$.
    \item[Mappatura delle regole di inferenza] 
        Per ogni asse del cubo, si dimostra che ogni regola di formazione dei tipi (sintattica) è isomorfa a una regola di introduzione o 
        eliminazione di un connettivo logico. L'albero di tipizzazione del termine $M$ riproduce esattamente i passi della derivazione logica.
    \item[Identità tra astrazione e scaricamento] 
        L'astrazione funzionale ($\lambda x:A$) mappa sul costrutto logico di \emph{scaricare un'ipotesi}. Il legame (\emph{binding}) sintattico 
        della variabile $x$ nel termine garantisce la chiusura della sezione corrispondente nell'albero di derivazione.
\end{description}

\paragraph{Corrispondenze dinamiche.}
\begin{description}[leftmargin=1.5em, style=nextline, font=\bfseries]
    \item[Riduzione come semplificazione] 
        La $\beta$-riduzione dei termini corrisponde al processo di \emph{cut-elimination} (o normalizzazione) della logica. 
        Calcolare un valore equivale a trasformare una dimostrazione indiretta in una \emph{forma normale} priva di passaggi 
        ridondanti.
    
    \item[Coerenza e non-abitabilità del falso] 
        La coerenza del sistema logico (l'impossibilità di dimostrare il falso $\bot$) mappa sulla non-abitabilità del tipo 
        vuoto. Un sistema di tipi solido (\emph{sound}) assicura che non esista alcun termine chiuso capace di testimoniare 
        una contraddizione logica.
\end{description}

\section{\texorpdfstring{$\lambda$}{Lambda}-calcolo semplicemente tipato con FADT}
Il $\lambda$-calcolo semplicemente tipato ($\lambda^{\to}$) presentato nel \ref{chap:type_systems}, pur essendo un ottimo modello di calcolo,
risualta limitato se osservato dal punto di vista del'isomorfismo di Curry-Howard. Esso infatti risulta isomorfo al solo \emph{frammento implicativo}
della logica intuizionista. Tuttavia una logica proposizionale completa richiede la capacita di esprimere congiunzioni ($\land$), disgiunzioni ($\lor$) 
e costanti di verità ($\top$) e falsità ($\bot$). Per sopperire a questa mancanza, è necessario arrichire il calcolo tramite i cosidetti \emph{Finite 
Algebraic Data Types (FADT)}. L'introduzione di tipi prodotto, somma, unit e void non risponde quindi a una mera esigenza di programmazione (come la 
definizione di strutture dati), ma alla necessità logica di fornire un'immagine computazionale a ogni connettivo della deduzione naturale. Tale sistema
viene indicato con $\lambda^{\to}_{\times, +, \mathbf{1}, \mathbf{0}}$.

\subsection{Grammatica}
La grammatica del sistema $\lambda^{\to}_{\times, +, \mathbf{1}, \mathbf{0}}$ estende il $\lambda^{\to}$ tramite un tipo per ogni costrutto della significativo
della logica proposizionale, con relatico costruttori e distruttori. Per costruttore si intende un costrutto sintattico che memorizza informazione e per distruttore
l'analogo che la estrae, similmente a come accade con l'astrazione e applicazione del $\lambda^{\to}$. 
\begin{equation}
    \begin{array}{lrll}
        \text{Types:}   & \tau, \sigma & ::= \iota                                      & \text{(Tipo base / atomico)}      \\
                        &              & \mid \tau \to \sigma                           & \text{(Tipo funzione)}            \\
                        &              & \mid \tau \times \sigma                        & \text{(Tipo prodotto)}            \\
                        &              & \mid \tau + \sigma                             & \text{(Tipo somma)}               \\
                        &              & \mid \mathbf{1}                                & \text{(Tipo Unit)}                \\
                        &              & \mid \mathbf{0}                                & \text{(Tipo Void)}                \\[6pt]
        \text{Terms:}   & e            & ::= x                                          & \text{(Variabile)}                \\
                        &              & \mid \lambda x : \tau . e                      & \text{(Astrazione)}               \\
                        &              & \mid e_1 \, e_2                                & \text{(Applicazione)}             \\
                        &              & \mid \langle e_1, e_2 \rangle                  & \text{(Coppia)}                   \\
                        &              & \mid \mathsf{fst} \ e                          & \text{(Proiezione sx)}            \\
                        &              & \mid \mathsf{snd} \ e                          & \text{(Proiezione dx)}            \\
                        &              & \mid \mathsf{inl} \ e                          & \text{(Iniezione sx)}             \\
                        &              & \mid \mathsf{inr} \ e                          & \text{(Iniezione dx)}             \\
                        &              & \mid \mathsf{case}\ e \ \mathsf{of} \ \dots    & \text{(Analisi per casi)}         \\
                        &              & \mid ()                                        & \text{(Costante Unit)}            \\
                        &              & \mid \mathsf{absurd}_\tau \ e                  & \text{(Eliminazione del falso)}       
    \end{array}
    \tag{STLC+}
    \label{eq:extended_stlc_grammar}
\end{equation}

\subsection{\texorpdfstring{$\beta$}{beta}-riduzione}
La semantica operazionale del sistema $\lambda^{\to}_{\times, +, \mathbf{1}, \mathbf{0}}$ definisce le regole di riduzione per i nuovi costrutti di dati. 
Ogni regola di $\beta$-riduzione specifica l'interazione tra un costruttore e il corrispondente distruttore, formalizzando il processo di calcolo per i tipi 
prodotto e somma.
\begin{equation}
    \begin{array}{lll}
        % Product redex
        \mathsf{fst} \, \langle e_1, e_2 \rangle & \to_\beta e_1 & \text{(Proiezione sx)} \\
        \mathsf{snd} \, \langle e_1, e_2 \rangle & \to_\beta e_2 & \text{(Proiezione dx)} \\[8pt]
        % Sum redex
        \mathsf{case} \, (\mathsf{inl} \, v) \, \mathsf{of} \, \{ \mathsf{inl} \, x \Rightarrow e_1 \mid \dots \} & \to_\beta e_1[v/x] & \text{(Ramo sx)} \\
        \mathsf{case} \, (\mathsf{inr} \, v) \, \mathsf{of} \, \{ \dots \mid \mathsf{inr} \, y \Rightarrow e_2 \} & \to_\beta e_2[v/y] & \text{(Ramo dx)}
    \end{array}
    \label{eq:fadt_beta_reduction}
\end{equation}

\paragraph{Nota sulla coerenza.} 
Si osservi che per il tipo $\mathbf{0}$ (falso) non esistono regole di riduzione. Ciò riflette la coerenza logica del sistema: in un contesto vuoto non 
è possibile costruire un termine di tipo $\mathbf{0}$, rendendo il costrutto $\mathsf{absurd}$ (ex falso quodlibet) un ramo di esecuzione logicamente 
irraggiungibile. Parallelamente, per il tipo $\mathbf{1}$ non sono necessarie riduzioni poiché non esiste un distruttore associato: la costante $()$ è 
già in forma normale.

\subsection{Type checking}
L'algoritmo di \emph{type checking} per il sistema $\lambda^{\to}_{\times, +, \mathbf{1}, \mathbf{0}}$ estende il processo definito per il $\lambda^{\to}$ integrando 
le regole di tipizzazione per i nuovi costruttori e distruttori.Mentre il costruttore della coppia ($\times$) richiede la presenza di entrambi i sottotermini per 
determinare il tipo in modo simmetrico, il caso dell'unione disgiunta ($+$) presenta una asimmetria: essendo presente un solo sottotermine ($\mathsf{inl}$ o $\mathsf{inr}$), 
è necessario esplicitare il tipo del ramo mancante come pedice del costruttore per preservare l'univocità del tipo.
\begin{equation}
    \text{type}(\Gamma, e) = 
    \begin{cases} 
        \dots               & \text{(Casi del $\lambda^{\to}$)}                                                                                     \\
        \tau \times \sigma  & \text{if } e = \langle e_1, e_2 \rangle \land \text{type}(\Gamma, e_1) = \tau \land \text{type}(\Gamma, e_2) = \sigma \\
        \tau                & \text{if } e = \mathsf{fst} \ e' \land \text{type}(\Gamma, e') = \tau \times \sigma                                   \\
        \sigma              & \text{if } e = \mathsf{snd} \ e' \land \text{type}(\Gamma, e') = \tau \times \sigma                                   \\
        \tau + \sigma       & \text{if } e = \mathsf{inl}_{\sigma} \ e' \land \text{type}(\Gamma, e') = \tau                                        \\
        \tau + \sigma       & \text{if } e = \mathsf{inr}_{\tau} \ e' \land \text{type}(\Gamma, e') = \sigma                                        \\
        \rho                & \text{if } 
            \begin{cases} 
                e = \mathsf{case} \ e_0 \ \mathsf{of} \ \mathsf{inl} \ x \Rightarrow e_1 \mid \mathsf{inr} \ y \Rightarrow e_2  \\
                \text{type}(\Gamma, e_0) = \tau + \sigma                                                                        \\
                \text{type}(\Gamma, x:\tau, e_1) = \rho                                                                         \\
                \text{type}(\Gamma, y:\sigma, e_2) = \rho 
            \end{cases}                                                                                                         \\
        \mathbf{1} & \text{if } e = ()                                                                      \\
        \tau       & \text{if } e = \mathsf{absurd}_{\tau} \ e' \land \text{type}(\Gamma, e') = \mathbf{0}  \\
        \perp      & \text{otherwise}
    \end{cases}
\end{equation}
Si osservi infine la natura delle regole riguardanti i tipi $\mathbf{1}$ e $\mathbf{0}$. Mentre il tipo \emph{Unit} è abitato dall'unica costante $()$, il tipo \emph{Void} 
non possiede costruttori espliciti: il termine $\mathsf{absurd}$ non produce un valore di tipo $\mathbf{0}$, ma ne richiede uno come input per derivare un tipo arbitrario $\tau$.

\subsubsection{Notazione relazionale}
L'algoritmo di \emph{type checking} $\mathsf{type}(\Gamma, e)$ appena presentato definisce un processo deterministico per verificare l'univocità dei tipi. Al fine di allineare questa 
trattazione alla formalizzazione logica necessaria per l'isomorfismo, da questo momento in poi esprimeremo l'output di tale algoritmo sotto forma di \emph{giudizio di tipizzazione} 
(o relazione di tipizzazione), utilizzando la seguente notazione standard del sistema formale.
\begin{equation}
    \Gamma \vdash e : \tau \iff \mathsf{type}(\Gamma, e) = \tau
\end{equation}
Il simbolo $\vdash$ (chiamato \emph{turnstile}) asserisce che, all'interno del contesto d'ipotesi $\Gamma$, il termine $e$ possiede validamente il tipo $\tau$. Questa scrittura ci permette di 
traslare la computazione ricorsiva dell'algoritmo in un sistema di regole d'inferenza statiche, dove la correttezza del programma coincide con l'esistenza di un \emph{albero di tipizzazione}. 
Tale formalismo relazionale è il prerequisito sintattico essenziale per enunciare i lemmi strutturali e mappare le regole informatiche su quelle della deduzione naturale.

\subsection{Lemmi fondamentali}
Al fine di poter enunciare e dimostrare formalemente l'isomorfismo di Curry-Howard tra la logica proposizionale intuizionista (IPC) e il sistema di tipi $\lambda^{\to}_{\times, +, \mathbf{1}, \mathbf{0}}$, 
è necessario stabilire alcuni lemmi fondamentali che garantiscano le proprietà essenziali che il sistema dei tipi deve soddisfare per essere considerato un modello fedele della logica. Tali lemmi riguardano 
la manipolazione dei termini e dei tipi all'interno del sistema, assicurando che le operazioni di sostituzione e indebolimento preservino la validità delle tipizzazioni.

\begin{lem}[Lemma d'indebolimento - weakening]
    Se $\Gamma \vdash e : \tau$ e $x \notin FV(e)$, allora per ogni tipo $\sigma$ si ha $\Gamma, x:\sigma \vdash e : \tau$.
\end{lem}
\begin{proof}
    La dimostrazione procede per induzione strutturale sul giudizio di tipizzazione o sulla sintassi del termine $e$. L'aggiunta di un'ipotesi inutilizzata non altera la 
    validità dei sotto-termini.
    \begin{description}[leftmargin=1.5em, style=nextline]
        \item [Caso Variabile ($e = y$)] 
            Se $\Gamma, y:\tau \vdash y : \tau$, allora anche $\Gamma, y:\tau, x:\sigma \vdash y : \tau$, poiché l'introduzione di $x$ non rimuove l'assegnamento preesistente di $y$ nel contesto.
        \item [Caso Astrazione ($e = \lambda y:\tau_1. e_1$)] 
            Per ipotesi induttiva su $e_1$, sappiamo che possiamo estendere il contesto. Assumendo per convenzione alfabetica che $y \neq x$ e $y \notin FV(e)$, se $\Gamma, y:\tau_1 \vdash e_1 : \tau_2$, 
            per ipotesi induttiva si ha $\Gamma, y:\tau_1, x:\sigma \vdash e_1 : \tau_2$. Applicando la regola di tipizzazione dell'astrazione, otteniamo $\Gamma, x:\sigma \vdash \lambda y:\tau_1. e_1 : \tau_1 \to \tau_2$.
        \item [Caso Coppia ($e = \langle e_1, e_2 \rangle$)] 
            Per le regole dei FADT, il tipo è $\tau_1 \times \tau_2$. Per ipotesi induttiva applicata sui singoli membri, se $\Gamma \vdash e_1 : \tau_1$ e $\Gamma \vdash e_2 : \tau_2$, allora si ha 
            $\Gamma, x:\sigma \vdash e_1 : \tau_1$ e $\Gamma, x:\sigma \vdash e_2 : \tau_2$. Riaplicando il costruttore prodotto si ottiene la tesi.
        \item [Caso Analisi ($e = \mathsf{case} \, e_0 \, \mathsf{of} \, \{ \mathsf{inl} \, y \Rightarrow e_1 \mid \mathsf{inr} \, z \Rightarrow e_2 \}$)] 
            L'indebolimento si distribuisce sul termine $e_0$ e, ipotizzando variabili le $y$ e $z$ non presenti del contesto, si applica induttivamente sui contesti estesi dei due rami 
            ($\Gamma, y:\tau_1, x:\sigma \vdash e_1 : \rho$) e ($\Gamma, z:\tau_2, x:\sigma \vdash e_2 : \rho$). L'unificazione dei rami preserva il tipo $\rho$.
        \item [Caso Unit ($e = ()$)] 
            Per le regole di tipizzazione, la costante () ha tipo $\mathbf{1}$, sotto qualsiasi contesto, senza alcuna dipendenza. Abbiamo quindi che $\Gamma \vdash () : \mathbf{1}$ implica 
            direttamente $\Gamma, x:\sigma \vdash () : \mathbf{1}$. Poiché per definizione $FV(()) = \emptyset$, l'estensione del contesto con l'ipotesi $x:\sigma$ non interferisce con la 
            validità del giudizio statico, soddisfacendo banalmente la tesi senza richiedere alcuna ipotesi induttiva.
        \item [Casi rimanenti ($\mathsf{fst}, \mathsf{snd}, \mathsf{inl}, \mathsf{inr}, \mathsf{absurd}$)]  
            Tali casi seguono per scomposizione induttiva diretta sui sotto-termini senza alterare la struttura del contesto.
    \end{description}
\end{proof}

\section{Dimostrazione del'isomorfismo IPC - \texorpdfstring{$\lambda^{\to}_{\times, +, \mathbf{1}, \mathbf{0}}$}{Lambda-Calculus with FADT}}
% TODO